\section{Overview}

FlexibleSUSY provides Mathematica and C++ code to create spectrum
generators for non-minimal supersymmetric models.

\subsection{Design goals}

FlexibleSUSY is designed with the following points in mind:

\paragraph{Speed}
Exploring the parameter space of supersymmetric models with a high
number of free parameters is quite time consuming.  Therefore
FlexibleSUSY aims to produce spectrum generators with a short run-time.

The two most time consuming parts of a spectrum generator are the
calculation of the beta functions and the pole masses for mixed
particles.

The reason is the following: The RG solving algorithms usually need
$O(10)$ iterations to calculate a particle spectrum with $0.1\%$
precision.  During each iteration the Runge--Kutta algorithm needs to
calculate all beta functions $O(50)$ times.  Most two-loop beta
functions involve $O(50)$ matrix multiplications and additions.  All
together one arrives at $O(2.5\cdot 10^4)$ matrix multiplications.  To
optimize these matrix multiplications FlexibleSUSY uses the fast
linear algebra package \href{Eigen}{http://eigen.tuxfamily.org}, which
exploits C++ expression templates to make the matrix multiplications
easy to write and optimize.

The second most time consuming part is the precise calculation of the
pole masses for mixed particles.  For each particle $\psi_i$ the full
self-energy matrix $\Sigma^\psi_{ij}(p=m^\text{tree}_{\psi_i})$ has to
be evaluated.  Each self-energy matrix entry again involves the
evaluation of $O(50)$ Feynman diagrams, which involves the calculation
of vertices and loop functions.  All in all, one arrives at $O(500)$
Feynman diagram or $O(3\cdot 10^4)$ loop function evaluations.  To
speed up the calculation of the pole masses FlexibleSUSY makes use of
multi-threading, where each pole mass is calculated in a separate
thread.  By this technique one can gain a speed-up of $20$--$30\%$.

\paragraph{Modularity}
The large variety of supersymmetric models makes it likely that the
user needs to modify the generated spectrum generator source code.
FlexibleSUSY uses object orientation with C++ to modularize the source
code to make it easy for the user to modify and extend the spectrum
generator.  An important application here are the constraints: The RG
solvers provide a constraint interface,
\figurename~\ref{fig:schematic-two-scale-constraint-interface},
%
\begin{figure}
  \centering
  \tikzumlset{fill class=white}
  \begin{tikzpicture}
    \umlclass[x=0, y=0]{RGFlow$<$Two\_scale$>$}{}{
      + solve()
    }
    \umlclass[x=6, y=0, type=abstract]{Constraint$<$Two\_scale$>$}{}{
      + \umlvirt{apply()}\\
      + \umlvirt{get\_scale() : double}
    }
    \umlclass[x=6, y=-3]{MyConstraint}{}{
      + apply()\\
      + get\_scale() : double
    }
    \umldep{RGFlow$<$Two\_scale$>$}{Constraint$<$Two\_scale$>$}
    \umlinherit{MyConstraint}{Constraint$<$Two\_scale$>$}
  \end{tikzpicture}
  \caption{Schematic two-scale constraint interface}
  \label{fig:schematic-two-scale-constraint-interface}
\end{figure}
%
which can be implemented to create the custom constraint
\code{MyConstraint}.  In the \code{apply()} function the constraint
application takes place and the \code{get_scale()} function specifies
the scale at which the constraint is to be applied.

\subsection{Requirements}

\begin{itemize}
\item Mathematica, version 7 or higher
\item SARAH, version 4.0.4 or higher \url{http://sarah.hepforge.org}
\item C++11 compatible compiler (g++ 4.4.7 or higher, clang++ 3.1 or
  higher)
\item Fortran compiler (gfortran, ifort etc.)
\item Eigen library, version 3.1 \url{http://eigen.tuxfamily.org}
\item Boost library, version 1.36.0 or higher
  \url{http://www.boost.org}
\item GNU scientific library \url{http://www.gnu.org/software/gsl/}
\item lapack (needed for the Lattice algorithm only)
  \url{http://www.netlib.org/lapack/}
\end{itemize}
%
Optional:
%
\begin{itemize}
\item Looptools, version 2.8 or higher
  \url{http://www.feynarts.de/looptools/}
\end{itemize}
